% META: {"id": "TEXP-MAT-EX-015", "chapitre": "TEXP-MATRICES-MARKOV", "type_objet": "exercice", "capacites": ["TEXP-MAT-C1"], "capacites_codes": ["C1"], "methodes": ["M1"], "parcours": 2, "competences": ["modeliser", "calculer", "raisonner"], "duree_min": 25, "mode_creation": "ex_nihilo", "sources_inspiration": [], "fichier_tex": "chapitres/TEXP-MATRICES-MARKOV/exercices/TEXP-MAT-EX-015.tex", "corrige_tex": "chapitres/TEXP-MATRICES-MARKOV/corriges/TEXP-MAT-CO-015.tex", "authoring": {"PEDAGOGICAL_GAP_ID": "GAP-E92D4ABEEFC9", "CAPACITY_ID": "C1", "EXERCISE_ROLE": "CONTEXT_VARIATION", "DIFFICULTY": 2, "AUTHORING_REASON": "les deux enonces existants donnent deux matrices sans provenance et demandent la meme liste d'operations ; ici les coefficients ont une unite, l'inverse sert a remonter d'une consommation a une production, et la derniere question demande si une operation licite a encore un sens", "ORIGINALITY_PROVENANCE": "ex_nihilo"}, "status": "generated"}
% BEGIN-VERIFY
% from sympy import Matrix, Rational, eye
% A = Matrix([[2, 5], [3, 4]])
% X = Matrix([[40], [30]])
% B = A * X
% assert B == Matrix([[230], [240]])
% assert A.det() == -7
% inverse = A.inv()
% assert inverse == Matrix([[Rational(-4, 7), Rational(5, 7)], [Rational(3, 7), Rational(-2, 7)]])
% assert A * inverse == eye(2)
% assert inverse * B == X
% Xprime = Matrix([[80], [30]])
% assert A * Xprime == Matrix([[310], [360]])
% D = Matrix([[0, 1], [1, 0]])
% Aprime = A + D
% assert Aprime == Matrix([[2, 6], [4, 4]])
% assert Aprime * X == Matrix([[260], [280]])
% assert A ** 2 == Matrix([[19, 30], [18, 31]])
% assert A * Matrix([[1], [0]]) == Matrix([[2], [3]])
% END-VERIFY
\begin{exercice}{TEXP-MAT-EX-015}{2}{25}
Un atelier fabrique deux modèles de vélo, $V_1$ et $V_2$. Chaque $V_1$ demande
$2$~heures-machine et $3$~kg d'aluminium ; chaque $V_2$ demande $5$~heures-machine
et $4$~kg d'aluminium. On note
\[
  A=\begin{pmatrix}2&5\\3&4\end{pmatrix},
  \qquad
  X=\begin{pmatrix}x_1\\x_2\end{pmatrix},
\]
où $x_1$ et $x_2$ sont les nombres de vélos de chaque modèle produits en une
semaine.
\begin{enumerate}
  \item Expliquer ce que représentent les deux coefficients de la matrice $AX$,
    en précisant leurs unités.
  \item Une semaine, l'atelier a consommé $230$~heures-machine et $240$~kg
    d'aluminium. Écrire l'équation matricielle correspondante.
  \item Calculer $\det(A)$, puis $A^{-1}$. Vérifier que $AA^{-1}=I_2$.
  \item Résoudre l'équation de la question~2 à l'aide de $A^{-1}$ et donner le
    nombre de vélos de chaque modèle.
  \item La semaine suivante, l'atelier double la production de $V_1$ sans changer
    celle de $V_2$. Calculer la nouvelle consommation.
  \item Un changement de fournisseur fait passer $V_2$ à $6$~heures-machine et
    $V_1$ à $4$~kg d'aluminium. Écrire la nouvelle matrice $A'$ comme une somme
    $A+D$, puis calculer la consommation correspondant à la production de la
    question~4.
  \item Le produit $A^2$ se calcule sans difficulté ; le faire. Ce résultat
    a-t-il une signification pour l'atelier ? Justifier en raisonnant sur les
    unités.
\end{enumerate}
\end{exercice}
